\documentclass[12pt]{article}
\usepackage[utf8]{inputenc}
\usepackage[T1]{fontenc}
\usepackage{hyperref}
\title{Kelo Language Resources for K-12 Curriculum Development}
\author{Compiled by Claude Sonnet 5.5 (claude-sonnet-5-5) for Andr\'as Kornai}
\date{}
\begin{document}
\maketitle

\section{Introduction}
Kelo (ISO 639-3 xel; Glottocode kelo1246) is a Nilo-Saharan language of the Eastern Jebel group. It is spoken by the Tornasi people in Blue Nile State, Sudan. Wikipedia gives about 200 native speakers (2009). UNESCO classifies it as critically endangered. Glottolog's assessment is ``endangered'', with low certainty. Searches found no teaching materials, primers, textbooks or school programmes for Kelo.

\section{Primary Educational Resources}
None were found. Kelo is not known to be taught in schools, and no native-language primers or readers turned up in search.

\section{Linguistic and Reference Materials}
\begin{itemize}
\item Wikipedia, Kelo language (overview, classification, speaker numbers): \url{https://en.wikipedia.org/wiki/Kelo_language}
\item Glottolog entry for Kelo: \url{https://glottolog.org/resource/languoid/id/kelo1246}
\item Catalogue of Endangered Languages entry: \url{https://www.endangeredlanguages.com/lang/3257?hl=en}
\end{itemize}
The search results mention M. Lionel Bender's 1997 article ``The Eastern Jebel Languages of Sudan'' (\emph{Afrika und \"Ubersee}) as the main descriptive source. No URL for it was seen, so it is not listed in the JSON.

\section{Religious and Community Resources}
Joshua Project has a people-group profile for Kelo/Tornasi in Sudan. It reports a mostly Muslim population: \url{https://joshuaproject.net/people_groups/12612}. No scripture or community-organisation resources in Kelo were found.

\section{Digital and Audio Resources}
No audio, video, corpus or software resources were found.

\section{Curriculum Development Context}
Documentation is minimal and the language is probably used only by the elderly. Any curriculum would first need new documentation work: recording elders, building a word list and agreeing an orthography. Sudan's current conditions make fieldwork difficult. The materials listed here are reference background only.
\end{document}